% =====================================================================
%  Section: Introduction
% =====================================================================
\section{Introduction}
\label{sec:intro}

Famous paintings encode not only colors and brushstrokes but also the cultural and
emotional intent of their creators. Bringing such static artworks to life---in the form
of short videos---has become a compelling research and educational objective, enabling
museum experiences, AI-driven storytelling, and creative content production. Recent
``living portrait'' demonstrations, in which the \textit{Mona Lisa} or Rembrandt
self-portraits are animated from a single image, illustrate the public appeal of this
direction~\cite{zakharov2019fewshot,guo2024liveportrait}.

Three technological threads have converged to make famous-painting video generation
possible: \textbf{(1) artistic style transfer}, which maps the visual appearance of a
reference painting onto a target image or video; \textbf{(2) image-to-video (I2V) and
video-to-video (V2V) generation}, which animate static content with temporally coherent
motion; and \textbf{(3) painting synthesis and stroke rendering}, which emulate the
physical creation process of paintings, including brushstroke dynamics. This survey
organizes representative works along these axes, with an emphasis on results published at
CVPR, ICCV, ECCV, AAAI, and IJCAI over the past decade.
